\src{176}
\uline{Soit} $\symup{\underline{\Omega}_X}$ (resp. $\symup{\underline{\Omega}_Y}$) \uline{le faisceau des germes de formes
différentielles régulières sur} $\symup{X}$ (resp. $\symup{Y}$). \uline{Alors on a des isomorphismes
canoniques} :
\begin{equation}\tag{3.6}
  \symup{\Extsh_{\Osh_X}^{p}\bigl(\underline{\Omega}_Y^{n-p},\underline{\Omega}_X^{n}\bigr)=\Osh_Y}
\end{equation}
\uline{ou encore}
\begin{equation}\tag{3.6 bis}
  \symup{\Extsh_{\Osh_X}^{p}\bigl(\Osh_Y,\underline{\Omega}_X^{n}\bigr)=\underline{\Omega}_Y^{n-p}}
\end{equation}
La formule (3.6 bis) peut servir de \uline{définition} de $\symup{\underline{\Omega}_Y^{n-p}}$ pour $\symup{Y}$
variété singulière. De façon précise :

\result{PROPOSITION}{5} \uline{Soit} $\symup{Y}$ \uline{un sous-ensemble algébrique de dimension}
$\symup{q=n-p}$ \uline{dans une variété algébrique non singulière} $\symup{X}$ \uline{de dimension} $\symup{n}$. \uline{Soient}
$\symup{F}$ \uline{un faisceau algébrique cohérent sur} $\symup{X}$ \uline{de support contenu dans} $\symup{Y}$, \uline{et} $\symup{L}$
\uline{un faisceau algébrique localement libre sur} $\symup{X}$. \uline{Alors les faisceaux}
$\symup{\Extsh_{\Osh_X}^{i}(F,L)}$ \uline{sont nuls pour} $\symup{i<p}$ ; \uline{tandis que pour} $\symup{i=p}$, \uline{on a un
isomorphisme canonique}
\begin{equation}\tag{3.7}
  \symup{\Extsh_{\Osh_X}^{p}(F,L)
  =\Homsh_{\Osh_X}\!\left(F,
  \Extsh^{p}(\Osh_X/J,L)\right)}
\end{equation}
\uline{où} $\symup{J}$ \uline{désigne un faisceau d'idéaux arbitraire sur} $\symup{X}$ \uline{annulant} $\symup{F}$ \uline{et dont
l'ensemble des zéros soit} $\symup{Y}$. \uline{En particulier, si} $\symup{F}$ \uline{est un faisceau algébrique
cohérent sur} $\symup{Y}$, \uline{on a}
\begin{equation}\tag{3.7 bis}
  \symup{\Extsh_{\Osh_X}^{p}(F,L)
  =\Homsh_{\Osh_Y}\!\left(F,\Extsh^{p}(\Osh_Y,L)\right).}
\end{equation}
\uline{Enfin}, $\symup{F}$ \uline{étant toujours un faisceau algébrique cohérent sur} $\symup{Y}$, \uline{les faisceaux}
$\symup{E^{i}(F)=\Extsh_{\Osh_X}^{p+i}(F,\underline{\Omega}_X^{n})}$ \uline{ne dépendent pas de la façon d'immerger l'espace algébrique} $\symup{Y}$ \uline{dans une variété
non singulière} $\symup{X}$.

La question étant locale, on peut supposer $\symup{X}$ affine et $\symup{L=\Osh_X}$. Cela ramène
à une question d'algèbre commutative, et même d'algèbre locale : Si $\symup{A}$ est une
localité régulière, $\symup{M}$ un $\symup{A}$-module dont le support est de dimension
$\symup{\leq q=n-p}$ prouver que $\symup{\operatorname{Ext}_{A}^{i}(M,A)=0}$ pour $\symup{i<p}$, et
$\symup{\operatorname{Ext}_{A}^{p}(M,A)
=\operatorname{Hom}_{A}\!\left(M,\operatorname{Ext}^{p}(A/J,A)\right)}$,
où $\symup{J}$ est un idéal quelconque de \fquote{dimension} $\symup{\leq q}$ annulant $\symup{M}$. Pour le
premier point, on procède par récurrence sur $\symup{q}$ : un dévissage immédiat ramène
au cas où $\symup{M}$ est de la forme $\symup{A/J}$, puis, en remplaçant $\symup{J}$ par un idéal plus
petit et utilisant l'hypothèse de récurrence, et la suite exacte des
$\symup{\operatorname{Ext}}$, au cas où $\symup{J}$ est engendré par un \fquote{système de paramètres}
comme dans la proposition 4 et où le résultat est immédiat. Le résultat
précédent implique que si $\symup{J}$ est un idéal

